\chapter*{Vorwort}
\addcontentsline{toc}{chapter}{Vorwort}

\begin{merksatz}
\boxtitel{Die Frage, die nach Band 23 bleibt}
Band 22 sagte vorher, was geschehen wird. Band 23 sagte, was sich ändert, wenn man handelt. Beides zusammen beantwortet die Frage des Arztes noch nicht.

\begin{center}
\begin{minipage}{0.9\textwidth}
\textbf{Denn zu wissen, was eine Behandlung bewirkt, heißt noch nicht zu wissen, ob man sie geben soll.}
\end{minipage}
\end{center}

\begin{center}
\begin{minipage}{0.88\textwidth}
\kursiv{Zwischen der Wirkung und der Entscheidung liegt eine Abwägung: \textbf{Die Behandlung verhindert Schlaganfälle und verursacht Blutungen.} Wie viele Blutungen sind ein verhinderter Schlaganfall wert? Diese Frage ist nicht empirisch --- keine Studie beantwortet sie. Sie ist auch nicht beliebig. Sie ist der Gegenstand dieses Bandes.}
\end{minipage}
\end{center}
\end{merksatz}

\begin{merksatz}
\boxtitel{Wo der Band steht}
\begin{center}
{\small
\renewcommand{\arraystretch}{1.52}
\begin{tabular}{@{}P{2.2cm}P{6.4cm}@{}}
\toprule
\textbf{Band 19} & Die Optimierung sucht das Beste --- bei bekannter Zielfunktion und bekannten Bedingungen.\\
\textbf{Band 21} & Die Spieltheorie fragt, was geschieht, wenn andere mitrechnen.\\
\textbf{Band 22} & Das Lernen sagt vorher.\\
\textbf{Band 23} & Der kausale Schluss sagt, was ein Eingriff bewirkt.\\
\textbf{Hier} & Wie wählt man, wenn man die Folgen nur mit Wahrscheinlichkeiten kennt und sie verschieden bewertet?\\
\bottomrule
\end{tabular}}
\end{center}

\begin{center}
\begin{minipage}{0.88\textwidth}
\kursiv{\textbf{Der Unterschied zu Band 19 ist die Unsicherheit, der zu Band 21 die Abwesenheit eines Gegenspielers.} Die Natur spielt nicht gegen uns; sie würfelt nur. Das macht die Sache einfacher --- und die Bewertung der Folgen schwieriger, denn sie muss von jemandem stammen.}
\end{minipage}
\end{center}
\end{merksatz}

\begin{merksatz}
\boxtitel{Drei Ergebnisse, die man nicht erwartet}
\begin{center}
{\small
\renewcommand{\arraystretch}{1.55}
\begin{tabular}{@{}P{1.0cm}P{7.4cm}@{}}
\toprule
$1$ & Wer nur wenige einleuchtende Regeln für vernünftige Vorlieben anerkennt, \kursiv{muss} einen Erwartungsnutzen maximieren --- eine Nutzenfunktion existiert dann, ob man will oder nicht. \kursiv{(Kapitel 2)}\\
$2$ & Ein Test lohnt sich nur in einem mittleren Bereich der Vortestwahrscheinlichkeit. Bei sehr niedriger \kursiv{und} bei hoher Wahrscheinlichkeit ist er überflüssig --- auch wenn er gut ist. \kursiv{(Kapitel 4)}\\
$3$ & Eine Gruppe vernünftiger Entscheider kann keine vernünftige gemeinsame Rangfolge bilden --- außer es entscheidet einer allein. \kursiv{(Kapitel 13)}\\
\bottomrule
\end{tabular}}
\end{center}

\begin{center}
\begin{minipage}{0.88\textwidth}
\kursiv{Das zweite Ergebnis widerspricht der klinischen Gewohnheit, im Zweifel zu testen. \textbf{Es gibt zwei Schwellen, nicht eine} --- und zwischen ihnen liegt der einzige Bereich, in dem ein Test die Entscheidung überhaupt ändern kann. Alles andere ist Aufwand ohne Folgen.}
\end{minipage}
\end{center}
\end{merksatz}

\begin{merksatz}
\boxtitel{Die sechs Teile}
\begin{center}
{\small
\renewcommand{\arraystretch}{1.5}
\begin{tabular}{@{}P{1.4cm}P{3.6cm}P{3.0cm}@{}}
\toprule
\textbf{Teil} & \textbf{Thema} & \textbf{Kapitel}\\
\midrule
I & Grundlagen & $1$--$3$\\
II & Klinische Entscheidungen & $4$--$6$\\
III & Über die Zeit & $7$--$9$\\
IV & Grenzen des Modells & $10$--$12$\\
V & Entscheiden für andere & $13$--$14$\\
VI & Stand & $15$\\
\bottomrule
\end{tabular}}
\end{center}

\schritt{Der Leitfall}
Frau B., $78$ Jahre, Vorhofflimmern. Ohne Gerinnungshemmung erleidet sie mit etwa sechs Prozent Wahrscheinlichkeit im Jahr einen Schlaganfall, mit Gerinnungshemmung mit zwei. Dafür steigt ihr Blutungsrisiko von einem auf drei Prozent.

\begin{center}
\begin{minipage}{0.88\textwidth}
\kursiv{\textbf{Vier verhinderte Schlaganfälle je tausend Behandlungsjahre, zwanzig zusätzliche Blutungen.} Ob das ein guter Tausch ist, hängt davon ab, wie schwer ein Schlaganfall gegenüber einer Blutung wiegt --- und das hängt von Frau B.\ ab. Dieser Fall trägt den ganzen Band: In Kapitel~1 wird er gerechnet, in Kapitel~4 mit einem Test versehen, in Kapitel~6 über zwanzig Jahre verfolgt und in Kapitel~14 gegen die Ansprüche anderer Patienten gestellt.}
\end{minipage}
\end{center}
\end{merksatz}

\begin{warnung}
\boxtitelr{Drei Warnungen}
\begin{center}
{\small
\renewcommand{\arraystretch}{1.55}
\begin{tabular}{@{}P{2.2cm}P{6.4cm}@{}}
\toprule
\textbf{Erste} & Die Rechnung liefert kein Ergebnis ohne Bewertung. Wer keine Nutzenwerte angibt, hat sie trotzdem benutzt --- unausgesprochen.\\
\textbf{Zweite} & Die Verfahren sind Hilfen zum Nachdenken, nicht Ersatz dafür. Wer ihnen blind folgt, hat die Verantwortung nur verschoben.\\
\textbf{Dritte} & Was für den einzelnen Patienten richtig ist, kann für die Gesamtheit falsch sein --- und umgekehrt.\\
\bottomrule
\end{tabular}}
\end{center}

\begin{center}
\begin{minipage}{0.88\textwidth}
\kursiv{Die erste Warnung ist die wichtigste. \textbf{Jede klinische Entscheidung enthält eine Abwägung}, auch die gewohnheitsmäßige; die Rechnung macht sie nur sichtbar. Das ist genau der Gewinn, den schon Band 23 für die kausalen Annahmen beanspruchte: nicht Gewissheit, sondern eine Form, in der man streiten kann.}
\end{minipage}
\end{center}
\end{warnung}

\begin{merksatz}
\boxtitel{Was dieser Band voraussetzt}
\begin{center}
{\small
\renewcommand{\arraystretch}{1.48}
\begin{tabular}{@{}P{3.0cm}P{1.6cm}P{3.0cm}@{}}
\toprule
\textbf{Werkzeug} & \textbf{Band} & \textbf{gebraucht in}\\
\midrule
Erwartungswert & $4$ & durchgehend\\
Satz von Bayes & $4$, $23$ & Kapitel 3--4\\
Optimierung & $19$ & Kapitel 7--9\\
Markow-Ketten & $4$, $16$ & Kapitel 6, 8\\
Spieltheorie & $21$ & Kapitel 13\\
\bottomrule
\end{tabular}}
\end{center}

\begin{center}
\begin{minipage}{0.88\textwidth}
\kursiv{\textbf{Teil~I und II kommen mit dem Erwartungswert und dem Satz von Bayes aus.} Wer Band 4 kennt, kann sofort beginnen. Die dynamische Programmierung aus Teil~III wird in Kapitel~7 von Grund auf entwickelt; Band 19 hilft, ist aber nicht nötig.}
\end{minipage}
\end{center}
\end{merksatz}
